% !TeX root = ../../../金融经济学_王江_讲义及习题.tex
% 来源：教材图2.1、2.2，印刷第10、12页；两状态下三维非负正交域。
\begin{center}
\begin{minipage}{\linewidth}
\WJFigureStart{fig:wj-endowment-consumption}%
\centering
\begin{tikzpicture}[>=Stealth,font=\small,x={(1cm,0cm)},y={(0cm,1cm)},z={(-.65cm,-.45cm)}]
 \draw[->] (0,0,0)--(2.7,0,0) node[right] {$e_{k,0}$};
 \draw[->] (0,0,0)--(0,2.5,0) node[above] {$e_{k,1a}$};
 \draw[->] (0,0,0)--(0,0,1.7) node[below left] {$e_{k,1b}$};
 \draw[gray,dashed] (1.8,1.8,1)--(1.8,0,1)--(1.8,0,0)
   (1.8,0,1)--(0,0,1) (1.8,1.8,1)--(0,1.8,0);
 \draw[->,blue!70!black,very thick] (0,0,0)--(1.8,1.8,1);
 \fill[blue!70!black] (1.8,1.8,1) circle(2pt)
   node[above right] {$e_k=\begin{bmatrix}e_{k,0}\\[1pt]\begin{bmatrix}e_{k,1a}\\e_{k,1b}\end{bmatrix}\end{bmatrix}$};
 \node at (1,-1.3,0) {给定的禀赋向量};
\end{tikzpicture}
\hspace{1cm}
\begin{tikzpicture}[>=Stealth,font=\small,x={(1cm,0cm)},y={(0cm,1cm)},z={(-.65cm,-.45cm)}]
 \draw[->] (0,0,0)--(2.7,0,0) node[right] {$c_{k,0}$};
 \draw[->] (0,0,0)--(0,2.5,0) node[above] {$c_{k,1a}$};
 \draw[->] (0,0,0)--(0,0,1.7) node[below left] {$c_{k,1b}$};
 \draw[gray,dashed] (1.8,1.8,1)--(1.8,0,1)--(1.8,0,0)
   (1.8,0,1)--(0,0,1) (1.8,1.8,1)--(0,1.8,0);
 \draw[->,blue!70!black,very thick] (0,0,0)--(1.8,1.8,1);
 \fill[blue!70!black] (1.8,1.8,1) circle(2pt)
   node[above right] {$c_k=\begin{bmatrix}c_{k,0}\\[1pt]\begin{bmatrix}c_{k,1a}\\c_{k,1b}\end{bmatrix}\end{bmatrix}$};
 \fill[black!45] (1.1,.8,.5) circle(1.5pt) (2,1.3,.3) circle(1.5pt);
 \node at (1,-1.3,0) {非负消费计划属于 $C$};
\end{tikzpicture}
\par\smallskip\small
两状态下，禀赋与消费计划都有三个分量：当期、未来状态 $a$、未来状态 $b$。消费集 $C=\mathrm R_+^3$ 包含整个非负正交域；预算约束将进一步限制可负担的消费。
\WJFigureCaption{禀赋向量与消费向量}
\end{minipage}
\end{center}
